% Adapted from academic-drawing/templates/tikz_template.tex to the book palette.
\begin{tikzpicture}[
  every node/.style={font=\figurefont\small,align=center,line width=.6pt},
  block/.style={book node,text width=39mm,minimum height=12mm,
    font=\figurefont\small},
  flow/.style={book arrow,rounded corners=1.2mm},
  train/.style={flow,dashed,draw=FigureLoss},
  feedback/.style={flow,dash dot,draw=BookAmber}
]
  \node[block,fill=FigureInput!12,draw=FigureInput] (material)
    at (-51mm,0) {输入材料\\来源、对象与时点};
  \node[block,fill=FigureInput!12,draw=FigureInput] (task)
    at (0,0) {任务要求\\语言、篇幅与格式};
  \node[block,fill=FigureLoss!10,draw=FigureLoss] (reference)
    at (51mm,0) {训练参考\\目标文本与结束符};
  \node[block,fill=FigureModel!12,draw=FigureModel,text width=62mm] (candidate)
    at (0,-25mm) {形成候选\\模板、选择、编辑或序列生成};
  \node[block,fill=FigureModel!8,draw=FigureModel] (translation)
    at (-51mm,-52mm) {翻译检查\\原意、实体与否定};
  \node[block,fill=FigureModel!8,draw=FigureModel] (summary)
    at (0,-52mm) {摘要检查\\选择、支持与遗漏};
  \node[block,fill=FigureModel!8,draw=FigureModel] (editing)
    at (51mm,-52mm) {编辑检查\\允许修改与意义保持};
  \node[block,fill=FigureOutput!12,draw=FigureOutput,text width=62mm] (output)
    at (0,-78mm) {通过：恢复并交付文本\\失败：保留错误与证据};

  \draw[flow] (material.south) |- (candidate.west);
  \draw[flow] (task.south) -- (candidate.north);
  \draw[train] (reference.south) |- node[pos=.25,right] {参数监督} (candidate.east);
  \coordinate (split) at (0,-39mm);
  \draw[flow] (candidate.south) -- (split) -- (summary.north);
  \draw[flow] (split) -| (translation.north);
  \draw[flow] (split) -| (editing.north);
  \coordinate (join) at (0,-66mm);
  \draw[flow] (translation.south) |- (join);
  \draw[flow] (editing.south) |- (join);
  \draw[flow] (summary.south) -- (output.north);
  \draw[feedback] (output.east) -- (79mm,-78mm) -- (79mm,-29mm)
    -- node[below,font=\figurefont\footnotesize] {重新选择或修订}
    ([yshift=-4mm]candidate.east);
\end{tikzpicture}
